\documentclass[11pt]{article}

\usepackage[margin=1in]{geometry}
\usepackage[T1]{fontenc}
\usepackage[utf8]{inputenc}
\usepackage{lmodern}
\usepackage{enumitem}
\usepackage{hyperref}
\usepackage{xcolor}

\hypersetup{
  colorlinks=true,
  linkcolor=blue,
  urlcolor=blue,
  citecolor=blue
}

\setlist[itemize]{leftmargin=*, itemsep=0.25em, topsep=0.25em}
\setlist[enumerate]{leftmargin=*, itemsep=0.25em, topsep=0.25em}
\setlength{\parindent}{0pt}
\setlength{\parskip}{0.55em}

\title{\vspace{-2em}Explanation of Revisions}
\author{}
\date{}

\begin{document}
\maketitle
\vspace{-3em}

This document describes the revisions made to the manuscript
\emph{A2M: Trace-Optimized Agent Hijacking in the MCP Ecosystem}
since its previous ARR submission.
The document is written to preserve anonymity
and is not a copy of the in-cycle author response.

\section{Overview of Major Revisions}

We made substantial revisions in response to the previous reviews
and the meta-review.
The main changes are summarized below.

\begin{itemize}
  \item \textbf{Revised novelty framing.}
  We removed or softened claims that MCP tool integration is a wholly novel
  attack surface.
  The paper now frames the contribution as a systematic method
  for exploiting a semantic supply-chain risk
  through coupled metadata optimization and trace-guided return-payload
  refinement.
  This change is reflected in the Abstract, Introduction,
  Related Work, and Threat Model sections.

  \item \textbf{Expanded related work.}
  We substantially revised the discussion of related MCP security work,
  including malicious MCP servers, parasitic toolchains,
  preference manipulation, metadata-based tool-selection attacks,
  and tool-selection perturbation attacks.
  The revised Related Work section now more clearly distinguishes A2M
  from prior metadata-only or tool-selection-focused attacks.

  \item \textbf{Clarified the two-stage formulation.}
  We moved the definitions of \emph{Attraction} and \emph{Manipulation}
  earlier in the paper and clarified that the two objectives are coupled
  during execution but are decoupled as optimization stages
  to make black-box search tractable.

  \item \textbf{Clarified the information exfiltration mechanism.}
  We revised the threat objective for Information Exfiltration
  to state that sensitive data is sent as arguments
  to the attacker-controlled malicious tool endpoint.
  This clarifies how exfiltration occurs without requiring attacker access
  to the victim's local agent logs.

  \item \textbf{Added and reorganized experiments.}
  We added AMA as a metadata-only attraction baseline,
  added cross-query generalization analysis,
  added failure-mode analysis,
  and added a ToolTweak-style paraphrase defense evaluation.
  To keep the main paper focused,
  detailed cross-query generalization and defense tables
  are reported in the appendix,
  with concise pointers in the main results section.

  \item \textbf{Expanded limitations and resource reporting.}
  We expanded the Limitations section to explicitly discuss
  system-prompt shifts, routing policies, tool renaming,
  safety wrappers, and changes in tool-pool composition.
  We also added aggregate token usage and estimated monetary cost
  for GLM-4.6 optimization and evaluation in the appendix.
\end{itemize}

\section{Response to the Meta-Review}

\textbf{Concern: Defense analysis.}
The meta-review noted that the previous version evaluated only
perplexity-based heuristics and a single auditor.
In Section 5.2 and Appendix E.1 (Defense Evaluation),
we added a ToolTweak-style paraphrase defense evaluation.
The revised paper now evaluates three defense families:
GPT-2 perplexity-based detection,
an LLM-based auditor implemented with Qwen3-8B,
and metadata paraphrasing.
The detailed defense results are reported in Tables 8 and 9,
with a concise pointer in the main results section.

\textbf{Concern: Robustness and generalization.}
In Section 5.2 and Appendix B.2 (Cross-Query Generalization),
we added a cross-query generalization study
to test whether optimized metadata overfits to a single query.
The revised appendix reports results in Table 5
on 50 held-out analogous queries
constructed from 10 base tasks.
A2M maintains high MTIR across all five evaluated models,
indicating that the Attraction stage captures reusable semantic cues
rather than only per-query artifacts.
We also expanded the Limitations section to explicitly acknowledge
remaining robustness dimensions not systematically varied in this paper,
including system prompts, routing policies, tool renaming,
safety wrappers, and tool-pool composition.

\textbf{Concern: Failure analysis.}
In Section 5.2 and Appendix D (Failure Analysis),
we added a failure-mode analysis.
The analysis groups unsuccessful attacks into three categories:
the malicious tool is not invoked,
the payload is ignored after invocation,
and the agent partially follows the payload
but fails to complete the objective.
Table 7 clarifies that, for IE and EIC,
ignored payloads are a dominant failure source,
which motivates runtime monitoring beyond selection-stage filtering.

\section{Response to Reviewer 1}

\subsection*{Concern 1: Novelty framing and attack-surface positioning}

The reviewer noted that prior work has already shown
that MCP tools, MCP servers, and tool outputs can be malicious.
We agree and revised the Abstract, Section 1, Section 2,
and Section 3.1 accordingly.
The Abstract no longer claims that MCP tool integration creates
a wholly novel attack surface.
Instead, it states that MCP systems expose
a semantic supply-chain risk
where attacker-controlled tool metadata and outputs
can influence tool selection and downstream reasoning.

In the Introduction and Threat Model,
we similarly replaced novelty-focused wording
with a more precise formulation:
A2M is presented as a systematic method
for exploiting the tool definition and return channels.
The Related Work section now explicitly discusses prior MCP security studies
on malicious servers, unsafe tool execution, parasitic toolchains,
preference manipulation, AMA, ToolHijacker, and ToolTweak.
The revised text distinguishes A2M by its full-chain optimization:
metadata optimization for selection
combined with trace-guided return-payload refinement after invocation.

\subsection*{Concern 2: Information exfiltration mechanism}

The reviewer asked how information is exfiltrated
if the attacker does not have direct access to victim-side logs.
In Section 3.3 (Attack Objectives),
we clarified the IE objective in the Threat Model.
The revised text states that the attacker manipulates the agent
into retrieving sensitive data from memory or the local file system
and sending it as arguments
to the attacker-controlled malicious tool endpoint.
Because the attacker controls the malicious MCP server endpoint,
the exfiltrated value is observable from the request arguments
received by that endpoint.

\subsection*{Concern 3: Baseline coverage}

In Section 5.1,
we revised the baseline discussion
to clarify that Zero-Shot Generation and LLM-GA
represent different attack paradigms:
single-pass unoptimized generation
and scalar-score black-box search without trace feedback.
We also added AMA as a metadata-only attraction baseline
and report the comparison in Section 5.2 and Table 2.
The revised results show that A2M's Attraction phase
achieves a higher average MTIR than AMA
($80.7$ versus $75.6$),
while acknowledging that AMA is stronger on some individual models.

\subsection*{Concern 4: Early definitions of Attraction and Manipulation}

In Section 1,
we moved the definitions of \emph{Attraction} and \emph{Manipulation}
to the Introduction.
The revised text defines Attraction as the selection-stage objective
of increasing malicious tool invocation probability,
and Manipulation as the post-invocation objective
of steering the agent's subsequent reasoning
toward an attacker-specified target.

\subsection*{Concern 5: Timeline and positioning of related works}

In Section 2,
we revised the Related Work section
to avoid describing relevant prior works as concurrent
when their release timelines indicate otherwise.
The revised section positions these works by technical focus
rather than by timeline:
metadata optimization and tool-selection bias
are discussed as primarily corresponding to the Attraction stage,
whereas A2M also optimizes return payloads
through trace-guided Manipulation.

\section{Response to Reviewer 2}

\subsection*{Concern 1: Cross-query generalization}

The reviewer noted that repeated per-query probing
may be unrealistic in deployed systems
and asked whether optimized metadata generalizes beyond the original query.
In Section 5.2 and Appendix B.2 (Cross-Query Generalization),
we added a cross-query generalization experiment.
We optimize malicious metadata on 10 base tasks
and evaluate the optimized metadata on 50 held-out analogous queries.
Table 5 reports MTIR ranging from $74.0$ to $90.0$
across the five evaluated models.
This suggests that the Attraction phase learns reusable semantic cues
rather than merely overfitting to one query.

\subsection*{Concern 2: Robustness under deployment drift}

In Sections 5.1 and 5.2,
we clarified the existing robustness evidence,
and in Limitations,
we expanded the discussion of deployment drift.
The main evaluation already spans 95 tasks,
six domains,
and toolsets drawn from 70 MCP servers,
which creates substantial variation in competing tools.
At the same time,
we now explicitly acknowledge that the current study
does not systematically vary system prompts,
routing policies,
tool renaming,
safety wrappers,
or the size and composition of competing tool pools.
These factors are now listed as future work in Limitations.

\subsection*{Concern 3: Baseline coverage}

In Section 5.1,
we added AMA as a representative metadata-only attraction baseline,
and in Section 5.2 and Table 2,
we report the comparison.
This complements the existing Zero-Shot and LLM-GA baselines
by directly evaluating tool-selection optimization.
The revised baseline section now describes each baseline
as representing a distinct attack paradigm.

\section{Response to Reviewer 3}

\subsection*{Concern 1: Relationship between Attraction and Manipulation}

The reviewer questioned whether Attraction and Manipulation
should be treated as unrelated goals,
since the quality of tool returns can affect subsequent behavior.
In Section 1,
we revised the Introduction to clarify that these objectives
are coupled during execution
but decoupled as optimization stages
to make the black-box search tractable.
This revision better reflects the role of the Analyzer--Optimizer loop:
after initial tool selection,
the return payload is refined using execution-trace feedback
to maintain agent trust and steer later reasoning.

\subsection*{Concern 2: Failure analysis for ASR results}

In Section 5.2 and Appendix D (Failure Analysis),
we added a failure analysis for ASR results.
For failed attacks on GLM-4.6,
we categorize failures as:
tool not invoked,
ignored payload,
and partial compliance.
Table 7 shows that ignored payloads are the dominant failure type
for IE and EIC,
while tool non-invocation and partial compliance
also explain failures in other scenarios.
This provides a clearer interpretation of the ASR results.

\subsection*{Concern 3: Defense evaluation}

In Section 5.2 and Appendix E.1 (Defense Evaluation),
we added a ToolTweak-style paraphrase defense evaluation
in addition to perplexity-based detection
and the LLM-based auditor.
Tables 8 and 9 report the detection and paraphrase-defense results.
The revised appendix shows that paraphrasing reduces MTIR
for some models,
but A2M still retains a $72.5$ average MTIR.
This supports the conclusion that metadata rewriting alone
does not reliably neutralize optimized attraction.

\subsection*{Concern 4: Figure 2 caption}

In Section 4,
we corrected the Figure 2 caption.
The revised caption removes the redundant ``Figure 2:'' prefix
and rewrites the caption for clarity.

\section{Additional Editorial and Organizational Changes}

We made several additional changes to improve clarity and readability.
We revised the ReAct and LangChain citations in the experimental setup.
We moved detailed cross-query generalization and defense evaluation tables
to the appendix,
while preserving concise pointers in the main results section.
We also added Appendix B.3 (Token Accounting) and Table 6,
which report aggregate token usage,
estimated GLM-4.6 cost under OpenRouter's Z.ai pricing,
and estimated average cost per task-scenario.

\end{document}
